\section*{Overview}

Centroid decomposition is the tree technique for problems where the relevant information depends on \emph{paths}, but a
full all-pairs or all-roots approach is too expensive. The decomposition repeatedly cuts the tree by balanced
separators, producing a recursion tree of depth \(O(\log n)\).

\section*{Problem-Driven Motivation}

Imagine a tree with \(n,q \le 2 \cdot 10^5\). Two online operations are supported:

\begin{itemize}
  \item mark one vertex red,
  \item query the distance from a vertex \(v\) to the nearest red vertex.
\end{itemize}

The naive answer for each query is a BFS or DFS from \(v\), which is \(O(n)\) per query. Even rerooting or one fixed
precomputation does not solve the dynamic update aspect cleanly.

What we need is a way to reduce every query and every update to only \(O(\log n)\) relevant tree nodes. That is exactly
what the centroid ancestor chain provides.

\section*{Recognition Pattern}

Centroid decomposition is a strong candidate when:

\begin{itemize}
  \item the graph is a tree,
  \item the property is about whole paths, not only subtrees,
  \item each query only needs path information aggregated through a few special separators,
  \item there are many updates and queries, and a full traversal per operation is too slow.
\end{itemize}

Typical contest signals are:

\begin{itemize}
  \item nearest active / colored vertex,
  \item count paths satisfying a distance condition,
  \item optimize over all paths through a chosen pivot,
  \item the statement is about a tree, but subtree-only tools like Euler tour or DSU on tree do not quite fit.
\end{itemize}

\section*{Derivation}

A centroid of a tree component is a vertex whose removal leaves no remaining component larger than half the size of the
current component.

Why is this useful?

\begin{itemize}
  \item If we recurse after removing the centroid, the recursion depth is \(O(\log n)\).
  \item Every original tree path belongs to exactly one of two categories at each recursive step:
    \begin{itemize}
      \item it passes through the current centroid,
      \item or it lies completely inside one recursive child component.
    \end{itemize}
\end{itemize}

So if we can process all paths that pass through the current centroid, recursion handles the rest.

\includegraphics[width=\textwidth]{figures/centroid-layers.svg}

For dynamic nearest-red queries, the derivation goes further:

\begin{itemize}
  \item every node stores the list of centroid ancestors above it,
  \item for each centroid \(c\), maintain the minimum distance from \(c\) to any red node,
  \item to query a node \(v\), walk through the centroid ancestors of \(v\) and combine
  \[
  dist(v,c) + best[c].
  \]
\end{itemize}

This is the entire trick. The balanced recursion makes the ancestor chain logarithmic.

\section*{Worked Problem}

\paragraph{Problem.}
Initially only node \(0\) is red. Support:

\begin{itemize}
  \item \texttt{paint(v)}: make \(v\) red,
  \item \texttt{nearest(v)}: return distance from \(v\) to the nearest red node.
\end{itemize}

\paragraph{Why naive fails.}
With \(2 \cdot 10^5\) operations, \(O(n)\) BFS per query is impossible.

\paragraph{Step 1: build the centroid tree.}
Every node gets a chain of centroid ancestors:
\[
v \rightarrow c_0 \rightarrow c_1 \rightarrow \cdots
\]
where \(c_0\) is the deepest centroid component containing \(v\), \(c_1\) is its parent in the centroid tree, and so
on.

\paragraph{Step 2: precompute distances to those ancestors.}
For every original node \(v\), store pairs \((c, dist(v,c))\) along that chain.

\paragraph{Step 3: process updates.}
When painting \(v\) red, update every centroid ancestor \(c\) of \(v\):
\[
best[c] = \min(best[c], dist(v,c)).
\]

\paragraph{Step 4: answer queries.}
For query node \(u\), try every centroid ancestor \(c\) of \(u\):
\[
answer = \min_c \bigl(best[c] + dist(u,c)\bigr).
\]

\paragraph{Why this works.}
Take the nearest red node \(r\) to \(u\). On the centroid decomposition recursion, there is a highest level where the
path \(u \leftrightarrow r\) is split by the current centroid \(c\). That path passes through \(c\), so
\[
dist(u,r) = dist(u,c) + dist(r,c).
\]
The update on \(r\) ensures that \(best[c] \le dist(r,c)\), so the query can recover the right answer from that
centroid level.

\section*{Implementation Reasoning}

The decomposition code has three distinct responsibilities:

\begin{itemize}
  \item compute subtree sizes inside the current component,
  \item find the centroid of that component,
  \item collect distances from every node in the component to this centroid and store them in the node's ancestor list.
\end{itemize}

The example implementation stores, for each original node, a vector of pairs:
\[
(\texttt{centroid}, \texttt{distance to centroid}).
\]
That is enough to support \(O(\log n)\) updates and queries for the nearest-red problem.

The important implementation invariant is:

\begin{quote}
For every node \(v\), its stored centroid path contains exactly the centroids of all decomposition components that
contain \(v\).
\end{quote}

\section*{Correctness Intuition}

Balanced separators are what make the method fast. The path-splitting argument is what makes it correct. Every query
finds the right separator level where the queried path is ``seen'' by one centroid, and every update has already pushed
its contribution into that centroid.

\section*{Complexity Analysis}

\begin{itemize}
  \item building the centroid decomposition: \(O(n \log n)\),
  \item one \texttt{paint}: \(O(\log n)\),
  \item one \texttt{nearest}: \(O(\log n)\),
  \item memory: \(O(n \log n)\) for storing centroid ancestor paths.
\end{itemize}

\section*{Common Pitfalls}

\begin{itemize}
  \item Confusing the original tree parent with the centroid-tree parent.
  \item Forgetting that the per-centroid processing still decides whether the final solution is fast enough.
  \item Computing distances to centroid ancestors lazily and accidentally turning queries back into \(O(n)\).
  \item Using centroid decomposition for a plain subtree problem that Euler tour or rerooting solves more directly.
\end{itemize}

\section*{Variants and Failure Modes}

\begin{itemize}
  \item Counting or optimizing over paths with a length threshold is a common offline centroid use.
  \item Weighted trees work too, but stored distances become weighted distances.
  \item If queries are about one root and one subtree only, heavy-light or Euler-tour reductions are usually cleaner.
  \item If the update/query structure is global but not path-based, centroid decomposition may be the wrong abstraction.
\end{itemize}

\section*{Practice Problems}

\begin{itemize}
  \item Nearest marked node on a tree.
  \item Count paths whose length satisfies a bound.
  \item Path-aggregation problems where a balanced separator naturally splits the path family.
\end{itemize}

\section*{References}

\begin{itemize}
  \item \href{https://cp-algorithms.com/graph/centroid_decomposition.html}{cp-algorithms: Centroid Decomposition}.
  \item \href{https://usaco.guide/plat/centroid?lang=cpp}{USACO Guide: Centroid Decomposition}.
  \item \href{https://codeforces.com/catalog}{Codeforces Catalog}.
\end{itemize}
